\section{Introduction}\label{sec:introduction}

Modern IoT and enterprise systems generate large volumes of interactions among devices, hosts, services, processes, and files. Detecting intrusions in these environments is difficult because malicious behavior is rare, evolves over time, and is only partially represented in labeled datasets. Signature-based systems depend on known attack patterns, while conventional statistical detectors often rely on fixed assumptions about normal and abnormal behavior. Both can struggle when behavior changes or an attack does not match a predefined pattern \cite{Buczak2016,sommer2010outside}. Machine learning (ML)-based intrusion detection is attractive because it can learn behavioral representations from data and use them to identify suspicious deviations \cite{Buczak2016}.

The same adaptability, however, is limited when ML models require comprehensive attack labels. Confirmed malicious samples are expensive to obtain, attacks are much rarer than normal activity, and previously collected labels can become less representative as attacker behavior changes \cite{sommer2010outside,caville2022anomale}. A supervised detector may therefore learn a strong boundary for known attacks yet fail when deployment behavior differs from its labeled training set \cite{sommer2010outside}. We consequently study a benign-only representation-learning setting: the detector learns the distribution of normal behavior and flags deviations as potential intrusions \cite{manzoor2016streamspot,caville2022anomale}. This setting does not guarantee detection of every unseen attack, but it reduces dependence on labeled attacks and is relevant to resilient detection under incomplete or changing threat knowledge.

System activity is inherently relational, which makes graphs a natural representation for intrusion detection. A provenance graph records how system entities and events depend on or interact with one another. Its nodes can represent hosts, devices, services, processes, or files, and its edges can represent communication, access, or system events. A single event may appear harmless in isolation while its role in a larger interaction pattern reveals malicious behavior. For example, port scanning can generate one-to-many connections \cite{lee2003detection}, lateral movement can create multi-step paths across hosts \cite{fang2022lmtracker,zhou2024lmdetect}, and coordinated attacks can concentrate activity around selected targets \cite{bakar2024ftgnet}. Graph-based detection can therefore evaluate events in context rather than treating each record independently \cite{manzoor2016streamspot,lo2022egraphsage}.

This relational context makes topology important because topology describes the pattern of connectivity and graph shape. Attacks that unfold through a system can leave paths, branches, mergers of connected components, or changes in connectivity across processes, files, services, and hosts \cite{fang2022lmtracker,zhou2024lmdetect}. Such signatures summarize how entities are connected rather than only which entities or events occur. They may therefore distinguish graphs that have similar local attributes but different global organization. A topology-aware representation complements local message passing by preserving structural information across multiple scales, which is useful when deployment-time anomalies are expressed as changes in interaction structure.

Graph contrastive learning provides a way to learn these representations without attack labels. It creates augmented views of the same graph and trains a graph neural network (GNN) encoder to recognize that the views originate from the same graph. GraphCL demonstrated that this strategy can learn graph-level representations without supervised labels \cite{you2020graphcl}. Yet standard GNN message passing primarily aggregates local neighborhoods. Benign and malicious graphs can share common local patterns while differing in security-relevant paths, branches, or connectivity. TopoGCL addresses this limitation by adding summaries of graph shape across multiple scales, making topology-aware contrastive learning a suitable starting point for graph intrusion detection \cite{chen2024topogcl}.

Despite this promise, TopoGCL was designed mainly for unsupervised graph classification, not for anomaly scoring in an intrusion detection system \cite{chen2024topogcl}. Classification and intrusion detection have different goals. In graph classification, labeled examples can help the model learn class boundaries after representation learning. In intrusion detection, there are no labeled attack examples during representation learning. The detector must decide whether a new graph is abnormal by comparing it with the learned distribution. Because of this, the embedding space used for distance-based scoring directly controls detection quality.

To address this gap, we propose \sysname, a topology-aware graph contrastive learning method for intrusion detection. \sysname learns from augmented views of benign graphs, combines message-passing and topology-aware embeddings, and scores each test graph by its distance from benign training embeddings. Crucially, scoring occurs in the projected space optimized during contrastive training, avoiding a mismatch between the learned objective and the representation used for detection. The method requires neither attack labels during representation learning nor an additional classifier.

We evaluate \sysname on StreamSpot, GraSec-IoT, and Unicorn Wget \cite{manzoor2016streamspot,hamouche2024graseciot,unicornwget}. StreamSpot contains heterogeneous system-activity graphs. GraSec-IoT contains communication and attack behavior from an Internet of Things setting. Unicorn Wget contains graph samples from a system-activity attack scenario. Together, these datasets test whether \sysname works across different graph construction settings and attack behaviors. We compare \sysname with a supervised GNN and GraphCL, a graph contrastive learning baseline.

In summary, this paper makes the following contributions:
\begin{itemize}
    \item We formulate graph intrusion detection as benign-only representation learning and anomaly scoring under limited attack-label availability.
    \item We introduce \sysname, which combines message-passing and topological embeddings and performs detection in the contrastive space optimized during training.
    \item We evaluate \sysname on three graph security datasets against supervised and benign-only baselines, finding higher F1 than GraphCL on StreamSpot and Unicorn Wget but lower F1 on GraSec-IoT.
\end{itemize}

Together, these contributions examine how topology-aware, benign-only anomaly detection can support resilient ML-driven security monitoring when labeled attacks are incomplete or may not represent deployment behavior.
